\textbf{Formation, persistence and the energetic value of paid adjacent fuel transport.}
\textbf{a}, The outer search combines four actual fixed two-kernel/response rules with two non-rigid initial density fields, a connected trilobe and a bridged long body. Initial mass is 71.2486104366 M and energy is 56.9988883493 EU; each resource stock initially totals 4096. All eight 128-update trajectories meet the frozen continuation criteria. Selection follows the observed formation processes and is fixed before stage 2, rather than ranking final mass. Six eligible bodies were not continued. Matrix labels abbreviate reference\_anisotropic, core\_ring, double\_shell\_far\_ring and compact\_ring\_core\_outer. Initial single-field windows are x=12--51, y=12--51, with no rotation or recentering. Rule and initial-shape effects are not separated by this search.

\textbf{b}, Each selected complete state at 128 is copied to transport-off and transport-on continuations; the only switched control is the external adjacent-fuel-transport capability. The reference\_anisotropic--bridged\_longbody candidate has become two separated blobs, not a connected long body. Both arms go exactly extinct at 220 and remain zero through 384. The core\_ring--connected\_trilobe candidate has become an elongated ring, not a persisting trilobe. Both arms remain alive and generate tissue through the observed endpoint 384. Displayed fields are archived transport-on density at 128/384 in the 64 domain and 384 in the embedded 96 domain for the ring and 128/192 for the blobs. The shared 128 field also represents the off start. The 64 ring windows are x=20--44, y=0--63; the 96 ring window is x=36--60, y=12--83. Grey ticks mark y=16 and y=80, the original resource-patch edges, not the new domain boundary. Blob windows x=0--63, y=24--43. Vertical coordinates increase downward. Cell sizes differ between the two morphological displays; each process keeps its own fixed window and scale. Single fields are nearest-neighbour images, linearly coloured from white at zero to teal at one M/cell, without contrast enhancement. The 64-by-64 domain is periodic. At saved state 232, density above 0.01 spans y=3--60, giving a seven-cell cross-boundary gap, whereas nonzero second-kernel support reaches vertical offsets of plus/minus eight. This is a geometric warning, not a calculation of the cross-boundary contribution or its causal onset. The reported 317 threshold contact is not the beginning of periodic effects; earlier energy differences at 192/256 are not asserted to be isolated from periodic geometry. The completed larger-domain comparison is included as described below; authorized P/F/P+F-transport followups are not included.

\textbf{c}, The six absolute paths (four original 64-domain arms and two new 96-domain ring arms) each use all 257 archived states from 128 through 384. For native GUI polylines, display geometry is simplified to within 0.20 screen pixels before pointer entry, retaining the shared start, extinction state 220, geometric-warning state 232, threshold-contact state 317 and endpoint; original arrays and CSVs remain unchanged. Both blob zero tails are drawn through 384. Within each domain, off curves are drawn wider than on curves to show their near overlap. Grey/teal identify off/on in 64; violet/orange identify off/on in 96. The two morphologies and both domains share the mass scale and share the energy scale, while M and EU are distinct units. The blob on-minus-off energy difference is temporarily negative (minimum -0.319840 EU at 150), later positive (maximum +3.751747 EU at 203), and zero after extinction; a zero endpoint does not imply identical paths. The ring endpoints are M=342.158713/343.119559 and E=352.220144/407.297181 for off/on.

\textbf{d}, All three relative changes share a percentage axis, computed against the ring's off arm within each domain. Teal and orange bars identify 64 and 96, respectively. Body-time is the archived sum of 0.05 times pre-update mass over actions 129--384: 4025.447864 versus 4033.817615 M-time. Final differences are +0.960846 M, +55.077037 EU and +8.369751 M-time. Fuel transfer acts along four-neighbour pairs supported by old tissue at both ends. Its actual edge flow totals 330.692931 F-edge and costs 0.02 EU per F-edge, paid once from the source's old energy, for 6.613859 EU. The transferred fuel is redistributed, not created. The same updated fuel field subsequently participates in uptake and generation. Total uptake is 403.340197/468.286488 EU and total F expenditure is 561.141113/627.616583. These recorded off-to-on totals are not a causal mediation decomposition. The blob on arm also pays 0.883558 EU for 44.177889 F-edge and increases uptake from 50.607189 to 56.968424 EU, without postponing extinction. Cumulative ring births are 375.716466/379.357369 M, while last-32 births are 21.376004/21.146327 M; cumulative improvement does not imply improvement in every late window.

Within-domain transport effects use complete shared receivers; cross-domain comparisons are conditional sensitivities of the embedding. All individual rules stay fixed, learning is disabled, and the shared actor continues its prescribed natural-expiry fresh-brake behavior; transport on/off is not an action-policy comparison. The selected receivers are not independent replicate histories. The first two jobs contain eight new 128-step first windows and four new 256-step continuations (2048 steps). The independent embedding job adds two 256-step continuations (512 steps), while its old-64 comparison reuses 512 archived updates already counted in the first source. Figure preparation adds no forward updates. Survival through 384 is a finite observation, not proof of permanent stability, resource autonomy, transport necessity or learned maintenance.


\textbf{Empty-periphery embedding.} The complete 128 state is embedded without changing its 18 stored arrays into coordinates 16:80 of a 96-by-96 periodic domain. The new periphery starts with all fields zero; total P/F and energy are conserved. Kernel integer offsets and values are retained, and the old central draw tape is reused while the peripheral tape is newly frozen and shared between the two arms. This changes domain geometry and numerical grid size; it does not identify a pure periodic-boundary mediator. Both new arms remain alive and produce tissue through 384 without threshold support touching the new outer boundary (minimum margin 15 cells for density above 0.01). The domain is still periodic, and the threshold omits small tails: this is not an infinite-domain isolation experiment.

In 96, off/on endpoints are M=357.410539/358.493166, E=375.325558/435.146177 EU and body-time=4188.174551/4196.680350 M-time. On-minus-off differences are +1.082627 M, +59.820619 EU and +8.505799 M-time. The actual on fee is 6.535624 EU for 326.781204 F-edge. Both arms lose mass at every update 297--384. Their last-32 births are 5.026233/4.859788 M and net mass changes are -7.226420/-7.083463 M; the corresponding old-64 births are 21.376004/21.146327 M with positive mass changes +2.197755/+2.317725 M. More endpoint tissue in 96 coexists with less late generation. The old-64 wrapped connection is not required for persistence over this observed new-domain window; neither domain establishes a stable or uniquely superior body.

At 96-domain state 384, positive requests remain 12.288018/12.319763 M per action, while admitted births are only 0.134469/0.122114 M. Ordered precursor rejections in the original core are 9.336942/10.245357 M; outer rejections are 0.885846/0.886583 M. These are stage-specific recorded quantities, not identified causal shares. Peripheral P and births are zero throughout; the final outer tissue (5.708047/5.711008 M) arrives through body transport rather than local generation. The P/F/P+F relocation experiment was only authorized when this figure was prepared and is not represented.
